\documentclass[12pt, a4paper]{article}

% ============================================================
%  Paquetes
% ============================================================
\usepackage[utf8]{inputenc}
\usepackage[T1]{fontenc}
\usepackage[spanish]{babel}
\usepackage{amsmath, amssymb}
\usepackage[ruled, vlined, linesnumbered, spanish]{algorithm2e}
\usepackage{geometry}
\usepackage{hyperref}
\usepackage{xcolor}
\usepackage{tikz}
\usetikzlibrary{positioning, arrows.meta, shapes.geometric}

\geometry{margin=2.5cm}

\title{Generación de Vecinos en MetaCARP\\
       \large Operadores de vecindario con ejemplos resueltos}
\author{}
\date{}

% ============================================================
\begin{document}
\maketitle
\tableofcontents
\newpage

% ------------------------------------------------------------
\section{Notación y representación}
% ------------------------------------------------------------

Una solución del CARP $s$ es un conjunto de $K$ rutas
$s = \{r_1, r_2, \ldots, r_K\}$.
Cada ruta $r_k = (e_1, e_2, \ldots, e_{n_k})$ es una secuencia ordenada de
aristas requeridas (tareas), identificadas por su etiqueta (p.\,ej.\
\texttt{TR1}, \texttt{TR2}, \dots).  Los tokens de depósito (\texttt{D}) se
eliminan antes de aplicar cualquier operador y se reinsertan para evaluar
el costo.

\begin{itemize}
    \item $r[i]$: tarea en la posición $i$ de la ruta $r$ (0-indexada).
    \item $r[i{:}j]$: subsecuencia desde la posición $i$ hasta $j$ inclusive.
    \item $r_a$, $r_b$: dos rutas distintas ($a \neq b$).
    \item $\mathit{head}(r, c) = r[0{:}c]$, $\mathit{tail}(r, c) = r[c{+}1{:}]$.
    \item $\|$ denota concatenación de listas.
\end{itemize}

% ------------------------------------------------------------
\section{Procedimiento general de generación de vecinos}
% ------------------------------------------------------------

El Algoritmo~\ref{alg:gen-neighbor} describe la rutina maestra
(\texttt{generar\_vecino} en \texttt{metacarp/vecindarios.py}) que toda
metaheurística llama en cada iteración.  Selecciona un operador mediante
ruleta ponderada (pesos provistos por \textsc{InterBias}) y reintenta hasta
500 veces si el operador elegido no es aplicable a la solución actual
(p.\,ej.\ \texttt{2opt\_intra} requiere una ruta con al menos 3 tareas).

Tres pasos rodean a la selección y aplicación del operador, y antes se
omitían en los pseudocódigos: la \textbf{normalización} que elimina el
token \texttt{D}, el \textbf{bucle de reintentos} (que hace \emph{continue}
en lugar de devolver \texttt{null}) y la \textbf{desnormalización} que
restaura \texttt{D} antes de devolver el vecino.

\begin{algorithm}[H]
\DontPrintSemicolon
\caption{\textsc{GenerarVecino} — Rutina maestra de generación de vecinos}
\label{alg:gen-neighbor}

\KwIn{Solución actual con depósito $s$;\; vector de pesos $w$ (o \texttt{null});\;
      conjunto de operadores $\mathcal{O}$;\; generador aleatorio $\mathit{rng}$;\;
      máximo de intentos $T_{\max}=500$}
\KwOut{Vecino con depósito $s'$;\; registro de movimiento $m$}

\BlankLine
$\hat{s} \leftarrow \textsc{NormalizarQuitarDepósito}(s)$
    \tcp{elimina el token "D" de inicio y fin de cada ruta}
$t \leftarrow 0$ \;

\Repeat{se encuentra un vecino válido}{
    $t \leftarrow t + 1$ \;
    \lIf{$t > T_{\max}$}{\textbf{lanzar} \texttt{RuntimeError} \tcp*{solución demasiado pequeña}}

    \BlankLine
    \eIf{$w \neq \texttt{null}$}{
        $\mathit{op} \leftarrow \textsc{SeleccionRuleta}(\mathcal{O},\, w,\, \mathit{rng})$
            \tcp{selección sesgada}
    }{
        $\mathit{op} \leftarrow \mathit{rng}.\textsc{choice}(\mathcal{O})$
            \tcp{selección uniforme}
    }

    \BlankLine
    $(\hat{s}',\, m) \leftarrow \textsc{AplicarOperador}(\mathit{op},\, \hat{s},\, \mathit{rng})$
        \tcp{si no es aplicable, \textbf{continue}}
}

$s' \leftarrow \textsc{RestaurarDepósito}(\hat{s}')$
    \tcp{añade "D" al inicio y fin de cada ruta antes de devolver}
\Return $(s',\, m)$
\end{algorithm}

\paragraph{Lectura autocontenida de los operadores.}
Por claridad, los pseudocódigos de cada operador en las
Secciones~\ref{sec:intra} y~\ref{sec:inter} se presentan
\emph{autocontenidos}: repiten la normalización, el bucle de reintentos, la
copia profunda y la desnormalización.  En el código real ese andamiaje vive
una sola vez en \textsc{GenerarVecino} (Algoritmo~\ref{alg:gen-neighbor}) y
cada operador aporta únicamente el filtrado de candidatas, el sorteo de
índices y el movimiento puro.

% ------------------------------------------------------------
\section{Operadores Intra-ruta}
\label{sec:intra}
% ------------------------------------------------------------

Los operadores intra-ruta actúan sobre una sola ruta $r_a$.
\textbf{Nunca cambian qué tareas pertenecen a qué ruta}, por lo que no
pueden reducir una violación de capacidad, pero son seguros de aplicar
libremente.

% ---- 3.1 Relocate-Intra ----
\subsection{\texttt{relocate\_intra}}

\textbf{Movimiento.} Quita la tarea en la posición $i$ de la ruta $r_a$ y
la reinserta en la posición $j \neq i$.

\textbf{Requisito.} $|r_a| \geq 2$.

\begin{algorithm}[H]
\DontPrintSemicolon
\caption{\texttt{relocate\_intra} --- Mover una tarea a otra posición dentro de la misma ruta}
\label{alg:relocate-intra}

\KwIn{Solución con depósito $s$;\; generador aleatorio $\mathit{rng}$;\; máximo de intentos $T_{\max}=500$}
\KwOut{Vecino con depósito $s'$;\; registro de movimiento $m$}

\BlankLine
$\hat{s} \leftarrow \textsc{NormalizarQuitarDepósito}(s)$
\tcp{elimina el token "D" de inicio y fin de cada ruta}
$t \leftarrow 0$\;
\Repeat{movimiento válido encontrado}{
    $t \leftarrow t + 1$\;
    \lIf{$t > T_{\max}$}{\textbf{lanzar} \texttt{RuntimeError}}
    $\mathcal{C} \leftarrow \{k : |\hat{r}_k| \geq 2\}$
    \tcp{rutas candidatas: al menos 2 tareas}
    \lIf{$\mathcal{C} = \emptyset$}{\textbf{continue}}
    $a \leftarrow \mathit{rng}.\textsc{choice}(\mathcal{C})$;\quad $n \leftarrow |\hat{r}_a|$\;
    $i \leftarrow \mathit{rng}.\textsc{randrange}(n)$
    \tcp{posición de origen, uniforme en $\{0,\ldots,n-1\}$}
    $j \leftarrow \mathit{rng}.\textsc{randrange}(n)$
    \tcp{posición de destino, uniforme en $\{0,\ldots,n-1\}$}
    \lIf{$i = j$}{\textbf{continue} \tcp*{no genera vecino distinto}}
    $\hat{s}' \leftarrow \textsc{CopiaProfunda}(\hat{s})$
    \tcp{Python pasa listas por referencia; sin copia se corrompería $\hat{s}$}
    $e \leftarrow \hat{r}'_a.\textsc{pop}(i)$\;
    $\hat{r}'_a.\textsc{insert}(j,\, e)$\;
    $m \leftarrow \textsc{MovimientoVecindario}(\texttt{relocate\_intra},\; \mathit{ruta\_a}=a,\; i=i,\; j=j)$\;
    $s' \leftarrow \textsc{RestaurarDepósito}(\hat{s}')$
    \tcp{añade "D" al inicio y fin de cada ruta antes de devolver}
    \Return $(s',\, m)$\;
}
\end{algorithm}

\noindent\textbf{Ejemplo visual.}
\[
    r_a = [\,\mathtt{TR1},\; \mathtt{TR2},\; \mathtt{TR3},\; \mathtt{TR4}\,]
    \xrightarrow{\;i=2,\; j=0\;}
    r'_a = [\,\mathtt{TR3},\; \mathtt{TR1},\; \mathtt{TR2},\; \mathtt{TR4}\,]
\]

% ---- 3.2 Swap-Intra ----
\subsection{\texttt{swap\_intra}}

\textbf{Movimiento.} Intercambia las tareas en las posiciones $i$ y $j$
($i \neq j$) dentro de la ruta $r_a$.

\textbf{Requisito.} $|r_a| \geq 2$.

\begin{algorithm}[H]
\DontPrintSemicolon
\caption{\texttt{swap\_intra} --- Intercambiar dos tareas dentro de la misma ruta}
\label{alg:swap-intra}

\KwIn{Solución con depósito $s$;\; generador aleatorio $\mathit{rng}$;\; máximo de intentos $T_{\max}=500$}
\KwOut{Vecino con depósito $s'$;\; registro de movimiento $m$}

\BlankLine
$\hat{s} \leftarrow \textsc{NormalizarQuitarDepósito}(s)$
\tcp{elimina el token "D" de inicio y fin de cada ruta}
$t \leftarrow 0$\;
\Repeat{movimiento válido encontrado}{
    $t \leftarrow t + 1$\;
    \lIf{$t > T_{\max}$}{\textbf{lanzar} \texttt{RuntimeError}}
    $\mathcal{C} \leftarrow \{k : |\hat{r}_k| \geq 2\}$
    \tcp{rutas candidatas: al menos 2 tareas}
    \lIf{$\mathcal{C} = \emptyset$}{\textbf{continue}}
    $a \leftarrow \mathit{rng}.\textsc{choice}(\mathcal{C})$;\quad $n \leftarrow |\hat{r}_a|$\;
    $(i,\, j) \leftarrow \mathit{rng}.\textsc{sample}(\{0,\ldots,n-1\},\, 2)$
    \tcp{dos índices distintos sin reemplazo (garantiza $i \neq j$ por construcción)}
    $\hat{s}' \leftarrow \textsc{CopiaProfunda}(\hat{s})$
    \tcp{Python pasa listas por referencia; sin copia se corrompería $\hat{s}$}
    $\hat{r}'_a[i],\, \hat{r}'_a[j] \leftarrow \hat{r}'_a[j],\, \hat{r}'_a[i]$
    \tcp{intercambio atómico: lado derecho se evalúa completo antes de asignar}
    $m \leftarrow \textsc{MovimientoVecindario}(\texttt{swap\_intra},\; \mathit{ruta\_a}=a,\; i=i,\; j=j)$\;
    $s' \leftarrow \textsc{RestaurarDepósito}(\hat{s}')$
    \tcp{añade "D" al inicio y fin de cada ruta antes de devolver}
    \Return $(s',\, m)$\;
}
\end{algorithm}

\noindent\textbf{Ejemplo visual.}
\[
    r_a = [\,\mathtt{TR1},\; \mathtt{TR2},\; \mathtt{TR3},\; \mathtt{TR4}\,]
    \xrightarrow{\;i=0,\; j=3\;}
    r'_a = [\,\mathtt{TR4},\; \mathtt{TR2},\; \mathtt{TR3},\; \mathtt{TR1}\,]
\]

% ---- 3.3 2opt-Intra ----
\subsection{\texttt{2opt\_intra}}

\textbf{Movimiento.} Invierte la subsecuencia $r_a[i{:}j]$ (ambos inclusive),
con $i < j$.

\textbf{Requisito.} $|r_a| \geq 3$ (segmento de al menos 2 tareas).

\begin{algorithm}[H]
\DontPrintSemicolon
\caption{\texttt{2opt\_intra} --- Invertir un segmento de tareas dentro de la misma ruta}
\label{alg:2opt-intra}

\KwIn{Solución con depósito $s$;\; generador aleatorio $\mathit{rng}$;\; máximo de intentos $T_{\max}=500$}
\KwOut{Vecino con depósito $s'$;\; registro de movimiento $m$}

\BlankLine
$\hat{s} \leftarrow \textsc{NormalizarQuitarDepósito}(s)$
\tcp{elimina el token "D" de inicio y fin de cada ruta}
$t \leftarrow 0$\;
\Repeat{movimiento válido encontrado}{
    $t \leftarrow t + 1$\;
    \lIf{$t > T_{\max}$}{\textbf{lanzar} \texttt{RuntimeError}}
    $\mathcal{C} \leftarrow \{k : |\hat{r}_k| \geq 3\}$
    \tcp{rutas candidatas: mínimo 3 tareas para garantizar segmento de longitud $\geq 2$}
    \lIf{$\mathcal{C} = \emptyset$}{\textbf{continue}}
    $a \leftarrow \mathit{rng}.\textsc{choice}(\mathcal{C})$;\quad $n \leftarrow |\hat{r}_a|$\;
    $i \leftarrow \mathit{rng}.\textsc{randrange}(0,\, n-1)$
    \tcp{inicio del segmento, uniforme en $\{0,\ldots,n-2\}$}
    $j \leftarrow \mathit{rng}.\textsc{randrange}(i+1,\, n)$
    \tcp{fin del segmento, uniforme en $\{i+1,\ldots,n-1\}$; garantiza $j > i$}
    $\hat{s}' \leftarrow \textsc{CopiaProfunda}(\hat{s})$
    \tcp{Python pasa listas por referencia; sin copia se corrompería $\hat{s}$}
    $\hat{r}'_a[i\,{:}\,j{+}1] \leftarrow \textsc{Revertir}(\hat{r}'_a[i\,{:}\,j{+}1])$
    \tcp{asignación de slice: reemplaza el segmento en su lugar}
    $m \leftarrow \textsc{MovimientoVecindario}(\texttt{2opt\_intra},\; \mathit{ruta\_a}=a,\; i=i,\; j=j)$\;
    $s' \leftarrow \textsc{RestaurarDepósito}(\hat{s}')$
    \tcp{añade "D" al inicio y fin de cada ruta antes de devolver}
    \Return $(s',\, m)$\;
}
\end{algorithm}

\noindent\textbf{Ejemplo visual.}
\[
    r_a = [\,\mathtt{TR1},\; \underbrace{\mathtt{TR2},\; \mathtt{TR3},\; \mathtt{TR4}}_{i=1,\;j=3},\; \mathtt{TR5}\,]
    \;\longrightarrow\;
    r'_a = [\,\mathtt{TR1},\; \mathtt{TR4},\; \mathtt{TR3},\; \mathtt{TR2},\; \mathtt{TR5}\,]
\]

% ------------------------------------------------------------
\section{Operadores Inter-ruta}
\label{sec:inter}
% ------------------------------------------------------------

Los operadores inter-ruta mueven tareas \textbf{entre dos rutas distintas}
$r_a$ y $r_b$ ($a \neq b$).  Pueden reducir una violación de capacidad en
una ruta sobrecargada transfiriendo demanda a otra ruta, por lo que se
prefieren (probabilidad $\alpha_{\mathrm{inter}} = 0.80$) cuando
\textsc{ViolacionCapacidad}$(s) > 0$.

% ---- 4.1 Relocate-Inter ----
\subsection{\texttt{relocate\_inter}}

\textbf{Movimiento.} Quita la tarea en la posición $i$ de la ruta $r_a$ e
inserta en la posición $j$ de la ruta $r_b$ ($a \neq b$).

\textbf{Requisito.} $|r_a| \geq 1$; $r_b$ puede estar vacía.

\begin{algorithm}[H]
\DontPrintSemicolon
\caption{\texttt{relocate\_inter} --- Mover una tarea de una ruta a otra ruta distinta}
\label{alg:relocate-inter}

\KwIn{Solución con depósito $s$;\; generador aleatorio $\mathit{rng}$;\; máximo de intentos $T_{\max}=500$}
\KwOut{Vecino con depósito $s'$;\; registro de movimiento $m$}

\BlankLine
$\hat{s} \leftarrow \textsc{NormalizarQuitarDepósito}(s)$
\tcp{elimina el token "D" de inicio y fin de cada ruta}
$t \leftarrow 0$\;
\Repeat{movimiento válido encontrado}{
    $t \leftarrow t + 1$\;
    \lIf{$t > T_{\max}$}{\textbf{lanzar} \texttt{RuntimeError}}
    $\mathcal{A} \leftarrow \{k : |\hat{r}_k| \geq 1\}$
    \tcp{rutas con al menos 1 tarea (candidatas como origen)}
    \lIf{$|\mathcal{A}| < 1$ \textbf{o} $|\hat{s}| < 2$}{\textbf{continue}}
    $\mathit{ra} \leftarrow \mathit{rng}.\textsc{choice}(\mathcal{A})$
    \tcp{ruta origen: uniforme entre rutas no vacías}
    $\mathit{rb} \leftarrow \mathit{rng}.\textsc{randrange}(|\hat{s}|)$
    \tcp{ruta destino: uniforme entre \textbf{todas} las rutas (puede estar vacía)}
    \lIf{$\mathit{ra} = \mathit{rb}$}{\textbf{continue} \tcp*{debe ser inter-ruta}}
    $i \leftarrow \mathit{rng}.\textsc{randrange}(|\hat{r}_{\mathit{ra}}|)$
    \tcp{posición en la ruta origen}
    $j \leftarrow \mathit{rng}.\textsc{randrange}(|\hat{r}_{\mathit{rb}}| + 1)$
    \tcp{posición en la ruta destino; $+1$ permite insertar al final}
    $\hat{s}' \leftarrow \textsc{CopiaProfunda}(\hat{s})$
    \tcp{Python pasa listas por referencia; sin copia se corrompería $\hat{s}$}
    $e \leftarrow \hat{r}'_{\mathit{ra}}.\textsc{pop}(i)$\;
    $\hat{r}'_{\mathit{rb}}.\textsc{insert}(j,\, e)$\;
    $m \leftarrow \textsc{MovimientoVecindario}(\texttt{relocate\_inter},\; \mathit{ruta\_a}=\mathit{ra},\; \mathit{ruta\_b}=\mathit{rb},\; i=i,\; j=j)$\;
    $s' \leftarrow \textsc{RestaurarDepósito}(\hat{s}')$
    \tcp{añade "D" al inicio y fin de cada ruta antes de devolver}
    \Return $(s',\, m)$\;
}
\end{algorithm}

\noindent\textbf{Ejemplo visual.}
\[
    \begin{array}{ll}
    r_a = [\mathtt{TR1},\; \mathtt{TR2},\; \mathtt{TR3}] & r_b = [\mathtt{TR4},\; \mathtt{TR5}]\\[4pt]
    \multicolumn{2}{c}{\xrightarrow{\;i=2\;(\mathtt{TR3}),\;\; j=1\;}}\\[4pt]
    r'_a = [\mathtt{TR1},\; \mathtt{TR2}] & r'_b = [\mathtt{TR4},\; \mathtt{TR3},\; \mathtt{TR5}]
    \end{array}
\]

% ---- 4.2 Swap-Inter ----
\subsection{\texttt{swap\_inter}}

\textbf{Movimiento.} Intercambia la tarea en la posición $i$ de la ruta
$r_a$ con la tarea en la posición $j$ de la ruta $r_b$.

\textbf{Requisito.} $|r_a| \geq 1$ y $|r_b| \geq 1$.

\begin{algorithm}[H]
\DontPrintSemicolon
\caption{\texttt{swap\_inter} --- Intercambiar una tarea entre dos rutas distintas}
\label{alg:swap-inter}

\KwIn{Solución con depósito $s$;\; generador aleatorio $\mathit{rng}$;\; máximo de intentos $T_{\max}=500$}
\KwOut{Vecino con depósito $s'$;\; registro de movimiento $m$}

\BlankLine
$\hat{s} \leftarrow \textsc{NormalizarQuitarDepósito}(s)$
\tcp{elimina el token "D" de inicio y fin de cada ruta}
$t \leftarrow 0$\;
\Repeat{movimiento válido encontrado}{
    $t \leftarrow t + 1$\;
    \lIf{$t > T_{\max}$}{\textbf{lanzar} \texttt{RuntimeError}}
    \lIf{$|\hat{s}| < 2$}{\textbf{continue}}
    $\mathcal{N} \leftarrow \{k : |\hat{r}_k| \geq 1\}$
    \tcp{rutas no vacías; ambas rutas participantes deben tener al menos 1 tarea}
    \lIf{$|\mathcal{N}| < 2$}{\textbf{continue}}
    $(\mathit{ra},\, \mathit{rb}) \leftarrow \mathit{rng}.\textsc{sample}(\mathcal{N},\, 2)$
    \tcp{dos rutas distintas no vacías, sin reemplazo; garantiza $\mathit{ra} \neq \mathit{rb}$}
    $i \leftarrow \mathit{rng}.\textsc{randrange}(|\hat{r}_{\mathit{ra}}|)$
    \tcp{posición en la ruta $\mathit{ra}$}
    $j \leftarrow \mathit{rng}.\textsc{randrange}(|\hat{r}_{\mathit{rb}}|)$
    \tcp{posición en la ruta $\mathit{rb}$}
    $\hat{s}' \leftarrow \textsc{CopiaProfunda}(\hat{s})$
    \tcp{Python pasa listas por referencia; sin copia se corrompería $\hat{s}$}
    $\hat{r}'_{\mathit{ra}}[i],\, \hat{r}'_{\mathit{rb}}[j] \leftarrow \hat{r}'_{\mathit{rb}}[j],\, \hat{r}'_{\mathit{ra}}[i]$
    \tcp{intercambio atómico: lado derecho se evalúa completo antes de asignar}
    $m \leftarrow \textsc{MovimientoVecindario}(\texttt{swap\_inter},\; \mathit{ruta\_a}=\mathit{ra},\; \mathit{ruta\_b}=\mathit{rb},\; i=i,\; j=j)$\;
    $s' \leftarrow \textsc{RestaurarDepósito}(\hat{s}')$
    \tcp{añade "D" al inicio y fin de cada ruta antes de devolver}
    \Return $(s',\, m)$\;
}
\end{algorithm}

\noindent\textbf{Ejemplo visual.}
\[
    \begin{array}{ll}
    r_a = [\mathtt{TR1},\; \mathtt{TR2}] & r_b = [\mathtt{TR3},\; \mathtt{TR4}]\\[4pt]
    \multicolumn{2}{c}{\xrightarrow{\;i=0\;(\mathtt{TR1}),\;\; j=1\;(\mathtt{TR4})\;}}\\[4pt]
    r'_a = [\mathtt{TR4},\; \mathtt{TR2}] & r'_b = [\mathtt{TR3},\; \mathtt{TR1}]
    \end{array}
\]

% ---- 4.3 2opt-Star ----
\subsection{\texttt{2opt\_star}}

\textbf{Movimiento.} Divide cada ruta en un punto de corte e intercambia
sus colas.  La ruta $r_a$ se divide en la posición $c_a$ en
$\mathit{head}_a = r_a[0{:}c_a]$ y $\mathit{tail}_a = r_a[c_a{+}1{:}]$.
Análogamente $r_b$ en $c_b$.  El resultado es:
\[
    r'_a = \mathit{head}_a \;\|\; \mathit{tail}_b,
    \qquad
    r'_b = \mathit{head}_b \;\|\; \mathit{tail}_a.
\]

\textbf{Requisito.} $|r_a| \geq 1$ y $|r_b| \geq 1$.

\begin{algorithm}[H]
\DontPrintSemicolon
\caption{\texttt{2opt\_star} --- Intercambiar las colas de dos rutas (2-opt*)}
\label{alg:2opt-star}

\KwIn{Solución con depósito $s$;\; generador aleatorio $\mathit{rng}$;\; máximo de intentos $T_{\max}=500$}
\KwOut{Vecino con depósito $s'$;\; registro de movimiento $m$}

\BlankLine
$\hat{s} \leftarrow \textsc{NormalizarQuitarDepósito}(s)$
\tcp{elimina el token "D" de inicio y fin de cada ruta}
$t \leftarrow 0$\;
\Repeat{movimiento válido encontrado}{
    $t \leftarrow t + 1$\;
    \lIf{$t > T_{\max}$}{\textbf{lanzar} \texttt{RuntimeError}}
    \lIf{$|\hat{s}| < 2$}{\textbf{continue}}
    $\mathcal{N} \leftarrow \{k : |\hat{r}_k| \geq 1\}$
    \tcp{rutas no vacías; ambas participantes necesitan al menos 1 tarea}
    \lIf{$|\mathcal{N}| < 2$}{\textbf{continue}}
    $(\mathit{ra},\, \mathit{rb}) \leftarrow \mathit{rng}.\textsc{sample}(\mathcal{N},\, 2)$
    \tcp{dos rutas distintas no vacías, sin reemplazo}
    $\mathit{cut\_a} \leftarrow \mathit{rng}.\textsc{randrange}(|\hat{r}_{\mathit{ra}}|)$
    \tcp{punto de corte en la ruta A, uniforme en $\{0,\ldots,|\hat{r}_{\mathit{ra}}|-1\}$}
    $\mathit{cut\_b} \leftarrow \mathit{rng}.\textsc{randrange}(|\hat{r}_{\mathit{rb}}|)$
    \tcp{punto de corte en la ruta B, uniforme en $\{0,\ldots,|\hat{r}_{\mathit{rb}}|-1\}$}
    $\hat{s}' \leftarrow \textsc{CopiaProfunda}(\hat{s})$
    \tcp{Python pasa listas por referencia; sin copia se corrompería $\hat{s}$}
    $\mathit{head\_A} \leftarrow \hat{r}'_{\mathit{ra}}[:\mathit{cut\_a}{+}1]$;\quad
    $\mathit{tail\_A} \leftarrow \hat{r}'_{\mathit{ra}}[\mathit{cut\_a}{+}1:]$\;
    $\mathit{head\_B} \leftarrow \hat{r}'_{\mathit{rb}}[:\mathit{cut\_b}{+}1]$;\quad
    $\mathit{tail\_B} \leftarrow \hat{r}'_{\mathit{rb}}[\mathit{cut\_b}{+}1:]$\;
    $\hat{r}'_{\mathit{ra}} \leftarrow \mathit{head\_A} \mathbin{\|} \mathit{tail\_B}$
    \tcp{ruta A adopta la cola de B}
    $\hat{r}'_{\mathit{rb}} \leftarrow \mathit{head\_B} \mathbin{\|} \mathit{tail\_A}$
    \tcp{ruta B adopta la cola de A}
    $m \leftarrow \textsc{MovimientoVecindario}(\texttt{2opt\_star},\; \mathit{ruta\_a}=\mathit{ra},\; \mathit{ruta\_b}=\mathit{rb},\; i=\mathit{cut\_a},\; j=\mathit{cut\_b})$\;
    $s' \leftarrow \textsc{RestaurarDepósito}(\hat{s}')$
    \tcp{añade "D" al inicio y fin de cada ruta antes de devolver}
    \Return $(s',\, m)$\;
}
\end{algorithm}

\noindent\textbf{Ejemplo visual} ($c_a=1$, $c_b=1$):
\[
    \begin{array}{ll}
    r_a = [\mathtt{TR1},\; \underbrace{\mathtt{TR2}}_{\text{head}},\; \overbrace{\mathtt{TR3},\;\mathtt{TR4}}^{\text{tail}}]
    & r_b = [\mathtt{TR5},\; \underbrace{\mathtt{TR6}}_{\text{head}},\; \overbrace{\mathtt{TR7}}^{\text{tail}}]\\[8pt]
    \multicolumn{2}{c}{\Downarrow}\\[4pt]
    r'_a = [\mathtt{TR1},\;\mathtt{TR2},\;\mathtt{TR7}]
    & r'_b = [\mathtt{TR5},\;\mathtt{TR6},\;\mathtt{TR3},\;\mathtt{TR4}]
    \end{array}
\]

% ---- 4.4 Cross-Exchange ----
\subsection{\texttt{cross\_exchange}}

\textbf{Movimiento.} Extrae un segmento contiguo $r_a[i{:}j]$ de $r_a$ y un
segmento contiguo $r_b[k{:}l]$ de $r_b$, y luego los intercambia:
\[
    r'_a = r_a[0{:}i] \;\|\; r_b[k{:}l] \;\|\; r_a[j{+}1{:}],
    \qquad
    r'_b = r_b[0{:}k] \;\|\; r_a[i{:}j] \;\|\; r_b[l{+}1{:}].
\]

\textbf{Requisito.} $|r_a| \geq 2$ y $|r_b| \geq 2$ (cada segmento debe
contener al menos dos tareas para que puedan cumplirse $j > i$ y $l > k$).

\begin{algorithm}[H]
\DontPrintSemicolon
\caption{\texttt{cross\_exchange} --- Intercambiar segmentos contiguos entre dos rutas}
\label{alg:cross-exchange}

\KwIn{Solución con depósito $s$;\; generador aleatorio $\mathit{rng}$;\; máximo de intentos $T_{\max}=500$}
\KwOut{Vecino con depósito $s'$;\; registro de movimiento $m$}

\BlankLine
$\hat{s} \leftarrow \textsc{NormalizarQuitarDepósito}(s)$
\tcp{elimina el token "D" de inicio y fin de cada ruta}
$t \leftarrow 0$\;
\Repeat{movimiento válido encontrado}{
    $t \leftarrow t + 1$\;
    \lIf{$t > T_{\max}$}{\textbf{lanzar} \texttt{RuntimeError}}
    \lIf{$|\hat{s}| < 2$}{\textbf{continue}}
    $\mathcal{N} \leftarrow \{k : |\hat{r}_k| \geq 2\}$
    \tcp{rutas con al menos 2 tareas en ambas participantes (requisito del segmento)}
    \lIf{$|\mathcal{N}| < 2$}{\textbf{continue}}
    $(\mathit{ra},\, \mathit{rb}) \leftarrow \mathit{rng}.\textsc{sample}(\mathcal{N},\, 2)$
    \tcp{dos rutas distintas con $\geq 2$ tareas, sin reemplazo}
    $n_a \leftarrow |\hat{r}_{\mathit{ra}}|$;\quad $n_b \leftarrow |\hat{r}_{\mathit{rb}}|$\;
    $i \leftarrow \mathit{rng}.\textsc{randrange}(0,\, n_a - 1)$
    \tcp{inicio del segmento en A, uniforme en $\{0,\ldots,n_a-2\}$}
    $j \leftarrow \mathit{rng}.\textsc{randrange}(i+1,\, n_a)$
    \tcp{fin del segmento en A, uniforme en $\{i+1,\ldots,n_a-1\}$; garantiza $j > i$}
    $k \leftarrow \mathit{rng}.\textsc{randrange}(0,\, n_b - 1)$
    \tcp{inicio del segmento en B, uniforme en $\{0,\ldots,n_b-2\}$}
    $l \leftarrow \mathit{rng}.\textsc{randrange}(k+1,\, n_b)$
    \tcp{fin del segmento en B, uniforme en $\{k+1,\ldots,n_b-1\}$; garantiza $l > k$}
    $\hat{s}' \leftarrow \textsc{CopiaProfunda}(\hat{s})$
    \tcp{Python pasa listas por referencia; sin copia se corrompería $\hat{s}$}
    $\mathit{seg\_A} \leftarrow \hat{r}'_{\mathit{ra}}[i\,{:}\,j{+}1]$;\quad
    $\mathit{seg\_B} \leftarrow \hat{r}'_{\mathit{rb}}[k\,{:}\,l{+}1]$\;
    $\hat{r}'_{\mathit{ra}} \leftarrow \hat{r}'_{\mathit{ra}}[{:}i] \mathbin{\|} \mathit{seg\_B} \mathbin{\|} \hat{r}'_{\mathit{ra}}[j{+}1{:}]$\;
    $\hat{r}'_{\mathit{rb}} \leftarrow \hat{r}'_{\mathit{rb}}[{:}k] \mathbin{\|} \mathit{seg\_A} \mathbin{\|} \hat{r}'_{\mathit{rb}}[l{+}1{:}]$\;
    $m \leftarrow \textsc{MovimientoVecindario}(\texttt{cross\_exchange},\; \mathit{ruta\_a}=\mathit{ra},\; \mathit{ruta\_b}=\mathit{rb},\; i=i,\; j=j,\; k=k,\; l=l)$\;
    $s' \leftarrow \textsc{RestaurarDepósito}(\hat{s}')$
    \tcp{añade "D" al inicio y fin de cada ruta antes de devolver}
    \Return $(s',\, m)$\;
}
\end{algorithm}

\noindent\textbf{Ejemplo visual} ($i=1, j=2, k=0, l=1$):
\[
    \begin{array}{ll}
    r_a = [\mathtt{TR1},\; \overbrace{\mathtt{TR2},\;\mathtt{TR3}}^{\mathit{seg}_a},\; \mathtt{TR4}]
    & r_b = [\overbrace{\mathtt{TR5},\;\mathtt{TR6}}^{\mathit{seg}_b},\; \mathtt{TR7}]\\[8pt]
    \multicolumn{2}{c}{\Downarrow}\\[4pt]
    r'_a = [\mathtt{TR1},\;\mathtt{TR5},\;\mathtt{TR6},\;\mathtt{TR4}]
    & r'_b = [\mathtt{TR2},\;\mathtt{TR3},\;\mathtt{TR7}]
    \end{array}
\]

% ------------------------------------------------------------
\section{Tabla Resumen}
% ------------------------------------------------------------

\begin{center}
\renewcommand{\arraystretch}{1.4}
\begin{tabular}{llccp{5cm}}
\hline
\textbf{Operador} & \textbf{Tipo} & \textbf{Tareas mín.} & \textbf{Rutas} & \textbf{Efecto} \\
\hline
\texttt{relocate\_intra} & intra & 2 & 1 & Mover una tarea a otra posición dentro de la misma ruta \\
\texttt{swap\_intra}     & intra & 2 & 1 & Intercambiar dos tareas dentro de la misma ruta \\
\texttt{2opt\_intra}     & intra & 3 & 1 & Invertir un segmento contiguo dentro de la misma ruta \\
\texttt{relocate\_inter} & inter & 1 & 2 & Mover una tarea de una ruta a otra \\
\texttt{swap\_inter}     & inter & 1+1 & 2 & Intercambiar una tarea entre dos rutas distintas \\
\texttt{2opt\_star}      & inter & 1+1 & 2 & Intercambiar las colas de dos rutas en puntos de corte elegidos \\
\texttt{cross\_exchange} & inter & 2+2 & 2 & Intercambiar segmentos contiguos entre dos rutas \\
\hline
\end{tabular}
\end{center}

\noindent\textbf{Nota.} Este documento cubre los siete operadores con
ejemplos resueltos.  Los operadores \texttt{or\_opt\_2} y \texttt{or\_opt\_3}
(reubicación de bloques de 2 y 3 tareas consecutivas) se documentan en la
sección completa de nueve operadores
(\texttt{pseudocodigos/operadores\_vecindario.tex}).

\end{document}
